% =============================================================================
% Relatório do TP de Teoria de Grafos e Computabilidade — PUC Minas 2026/1
%
% INSTRUÇÕES PARA USAR O TEMPLATE OFICIAL DA SBC:
%   1. Acesse: https://www.overleaf.com/latex/templates/sbc-conferences-template
%   2. Baixe os arquivos: sbc-template.cls e sbc.bst
%   3. Coloque-os nesta pasta (relatorio/) ao lado deste .tex
%   4. Substitua a linha \documentclass{article} abaixo por:
%        \documentclass[12pt]{sbc-template}
%      e remova os pacotes que o sbc-template já carrega (geometry, fontenc...)
%
% Este esqueleto está em \documentclass{article} para compilar imediatamente
% mesmo antes de você obter o template SBC, facilitando a iteração local.
% =============================================================================

\documentclass[a4paper,11pt]{article}

\usepackage[utf8]{inputenc}
\usepackage[T1]{fontenc}
\usepackage[brazil]{babel}
\usepackage[margin=2.5cm]{geometry}
\usepackage{graphicx}
\usepackage{listings}
\usepackage{xcolor}
\usepackage{hyperref}
\usepackage{amsmath,amssymb}
\usepackage{tikz}
\usetikzlibrary{positioning,arrows.meta}
\usepackage{booktabs}
\usepackage{caption}

\lstset{
  basicstyle=\ttfamily\small,
  keywordstyle=\color{blue!70!black}\bfseries,
  commentstyle=\color{gray},
  stringstyle=\color{green!50!black},
  showstringspaces=false,
  breaklines=true,
  numbers=left,
  numberstyle=\tiny\color{gray},
  frame=single,
  language=Python
}

\title{
  Análise da Rede de Colaboração em Repositórios do GitHub\\
  utilizando Teoria de Grafos
}
\author{
  % TODO: substituir pelos integrantes do grupo (nome completo e matrícula)
  Integrante 1 \and
  Integrante 2 \and
  Integrante 3 \and
  Integrante 4 \\[1ex]
  \small Pontifícia Universidade Católica de Minas Gerais\\
  \small Curso de Engenharia de Software\\
  \small Disciplina: Teoria de Grafos e Computabilidade — Prof. Leonardo V.\ Cardoso
}
\date{2026/1}

\begin{document}

\maketitle

\begin{abstract}
% TODO: resumir em 5--10 linhas o problema, a abordagem e os principais
% resultados obtidos. Falar do repositório escolhido, dos 4 grafos modelados,
% das métricas calculadas e do principal insight observado.
\end{abstract}

% -----------------------------------------------------------------------------
\section{Introdução}\label{sec:intro}

% Contextualize o problema: por que estudar colaboração em projetos
% open-source com grafos? Cite a motivação acadêmica (disciplina) e prática
% (entender comunidades de software). Apresente o repositório escolhido
% (>5.000 estrelas) e justifique a escolha (volume de dados, diversidade
% de interações, relevância da comunidade).
%
% Estrutura sugerida:
%   - Contexto e motivação
%   - Repositório escolhido (TODO: definir e justificar)
%   - Objetivos específicos do trabalho
%   - Organização do restante do relatório

% -----------------------------------------------------------------------------
\section{Modelagem do Problema}\label{sec:modelagem}

% Definição formal dos grafos. Reproduza aqui o quadro dos 4 grafos:
%   Grafo 1 (comentários), Grafo 2 (fechamento), Grafo 3 (reviews/merges),
%   Grafo Integrado (ponderado).
%
% Apresente a tabela de pesos exigida pelo enunciado:

\begin{table}[h]
  \centering
  \caption{Pesos das interações no grafo integrado (Etapa 1 do TP).}
  \label{tab:pesos}
  \begin{tabular}{lc}
    \toprule
    \textbf{Tipo de interação} & \textbf{Peso} \\
    \midrule
    Comentário em issue ou pull request           & 2 \\
    Fechamento de issue por outro usuário         & 3 \\
    Revisão / aprovação de pull request           & 4 \\
    Merge de pull request                         & 5 \\
    \bottomrule
  \end{tabular}
\end{table}

% Discuta: por que esses pesos refletem a intensidade da colaboração?
% Mencione que os pesos podem ser adaptados (enunciado permite) e justifique
% qualquer adaptação que vocês decidirem fazer.

% -----------------------------------------------------------------------------
\section{Arquitetura da Ferramenta}\label{sec:arquitetura}

% Apresente o diagrama de classes principal. Sugestão usando TikZ:
% (substitua/melhore conforme necessário)

\begin{figure}[h]
  \centering
  \begin{tikzpicture}[
    node distance=1.2cm and 2cm,
    every node/.style={draw, rounded corners, font=\small, align=center},
    abstract/.style={fill=blue!8},
    concrete/.style={fill=green!10},
    arrow/.style={-{Stealth[length=2mm]}, thick}
  ]
    \node[abstract] (abs) {AbstractGraph\\\textit{(abstrata)}};
    \node[concrete, below left=of abs] (lst) {AdjacencyListGraph};
    \node[concrete, below right=of abs] (mat) {AdjacencyMatrixGraph};
    \draw[arrow] (lst) -- (abs);
    \draw[arrow] (mat) -- (abs);
  \end{tikzpicture}
  \caption{Hierarquia da API de grafos.}
  \label{fig:classes}
\end{figure}

% Descreva os módulos:
%   - codigos/core/graph/     -> AbstractGraph e implementações
%   - codigos/github_miner/   -> coleta via API REST do GitHub
%   - codigos/services/       -> GraphBuilderService e MetricsService
%   - codigos/exporters/      -> GEXF, GraphML, CSV
%   - codigos/app/            -> CLI principal e api_demo
%
% Justifique:
%   - Por que duas implementações (lista vs matriz)? -> trade-off espaço/tempo
%   - Onde a herança (AbstractGraph) é usada e por quê
%   - Restrições atendidas: idempotência de addEdge, rejeição de self-loops,
%     exceções para índices inválidos, ausência de bibliotecas externas de
%     grafo (NetworkX, igraph etc.).

% -----------------------------------------------------------------------------
\section{Coleta e Processamento dos Dados}\label{sec:coleta}

% Descreva o fluxo do github_miner:
%   1. GithubClient autentica e consome /repos/{owner}/{repo}/...
%   2. RateLimiter respeita os 5000 req/h da API
%   3. GithubParser transforma JSON em modelos tipados (User, Issue, PR,
%      Review, Comment) com closed_by e merged_by capturados
%   4. DataTransformer agrega tudo em CollaborationGraph
%   5. GraphBuilderService gera os 4 grafos (Tabela~\ref{tab:pesos})

% Cite volume de dados coletados (TODO: rodar e preencher)
%   - X usuários únicos
%   - Y issues / Z pull requests
%   - W comentários
%   - K reviews / merges

% -----------------------------------------------------------------------------
\section{Métricas e Análise}\label{sec:metricas}

% Liste e explique cada métrica implementada (todas SEM bibliotecas externas
% de grafo, em \texttt{codigos/services/metrics.py}):

\subsection{Centralidade}
% Degree centrality, in/out-degree, betweenness (Brandes), closeness,
% PageRank (power iteration, damping 0.85), eigenvector (power iteration
% sobre A^T). Inclua a fórmula de cada uma.

\subsection{Estrutura e coesão}
% Densidade da rede, coeficiente de clustering local, assortatividade
% (correlação de Pearson sobre graus dos extremos).

\subsection{Comunidades}
% Detecção por Label Propagation Algorithm (LPA), modularidade Q
% (Newman, ponderada), bridging ties (arestas entre comunidades distintas
% ordenadas por peso).

% -----------------------------------------------------------------------------
\section{Resultados}\label{sec:resultados}

% TODO: rodar a CLI sobre o repositório escolhido e preencher.
% Inclua:
%   - Tabela de top-10 colaboradores por PageRank
%   - Tabela de top-10 por betweenness
%   - Gráfico de distribuição de graus
%   - Screenshot do grafo no Gephi (importar .gexf gerado)
%   - Modularidade Q obtida e número de comunidades detectadas
%   - Exemplos de bridging ties

\begin{figure}[h]
  \centering
  \fbox{\parbox{0.7\linewidth}{\centering
    \vspace{4cm}
    [Substituir por screenshot do grafo no Gephi]
    \vspace{4cm}
  }}
  \caption{Visualização do grafo integrado no Gephi.}
  \label{fig:gephi}
\end{figure}

% -----------------------------------------------------------------------------
\section{Responsabilidades dos Integrantes}\label{sec:responsabilidades}

% Exigência do enunciado: indicar o que cada integrante fez.
\begin{itemize}
  \item \textbf{Integrante 1}: % TODO
  \item \textbf{Integrante 2}: % TODO
  \item \textbf{Integrante 3}: % TODO
  \item \textbf{Integrante 4}: % TODO
\end{itemize}

% -----------------------------------------------------------------------------
\section{Conclusão}\label{sec:conclusao}

% Resuma:
%   - O que foi entregue (4 grafos, 12+ métricas, exportação Gephi, testes)
%   - Principais aprendizados sobre a rede analisada
%   - Limitações encontradas (rate limit da API, escolhas de modelagem)
%   - Trabalhos futuros (Louvain, métricas temporais, comparação entre repos)

% -----------------------------------------------------------------------------
\begin{thebibliography}{9}

\bibitem{brandes2001}
  Ulrik Brandes. \emph{A Faster Algorithm for Betweenness Centrality}.
  Journal of Mathematical Sociology, 25(2):163--177, 2001.

\bibitem{page1999}
  Lawrence Page, Sergey Brin, Rajeev Motwani, Terry Winograd.
  \emph{The PageRank Citation Ranking: Bringing Order to the Web}.
  Stanford InfoLab, 1999.

\bibitem{newman2010}
  M.\ E.\ J.\ Newman. \emph{Networks: An Introduction}.
  Oxford University Press, 2010.

\bibitem{raghavan2007}
  U.\ N.\ Raghavan, R.\ Albert, S.\ Kumara.
  \emph{Near linear time algorithm to detect community structures in
  large-scale networks}. Physical Review E, 76(3), 2007.

\bibitem{gephi}
  Mathieu Bastian, Sebastien Heymann, Mathieu Jacomy.
  \emph{Gephi: an open source software for exploring and manipulating networks}.
  ICWSM, 2009.

\end{thebibliography}

\end{document}
